\section{问题四的模型建立与求解}

\subsection{第四问改变了哪些条件}

问题四在问题三的基础上引入定向干扰源：源总数仍为 10 至 16 个且未知，但现在既有全向源又有定向源，定向源的个数与朝向也未知。由附录一，定向源的有效覆盖角度为发射方向两侧各 $90^\circ$（含边界），范围之外没有信号辐射；而光学精确定位与清除只受 20 米距离约束，与源的朝向无关。

以 $g$ 表示源位置、$R$ 表示接收半径、$n=(\cos\varphi,\sin\varphi)$ 表示定向源的发射方向单位向量、$\tau\in\{\text{全向},\text{定向}\}$ 表示类型，则位置 $s$ 能接收到某个未清除源，当且仅当
\begin{equation}
  \lVert s-g\rVert\le R
  \quad\text{且}\quad
  \bigl[\tau=\text{全向}\ \text{或}\ n^{T}(s-g)\ge0\bigr].
  \label{eq:directional-reception}
\end{equation}
需特别注意：示向度是 $s$ 指向 $g$ 的角度，\textbf{不是}发射方向 $\varphi$。因此由示向度不能推断源的类型或朝向；由无信号也不能认定附近没有源。

据此逐项检查第三问的机制，得到表 \ref{tab:q4-inherit} 的继承清单。这张清单是第四问建模的起点：真正需要重新建立的只有"无信号的含义"与"覆盖终止证书"两项，其余几何与清除模块可以原样继承。

\begin{table}[H]
  \centering
  \caption{第三问机制在第四问中的继承情况}
  \label{tab:q4-inherit}
  \begingroup
  \renewcommand{\arraystretch}{1.4}
  \begin{tabularx}{\textwidth}{L{0.36\textwidth}Y}
    \toprule[1.2pt]
    {\heifont\bfseries 第三问机制} & {\heifont\bfseries 第四问的处理} \\
    \midrule[0.5pt]
    有界示向度半平面、多边形外包、最小包围圆 & 收到信号时继续成立，原样继承 \\
    最近保证清除点计算 & 光学不受朝向限制，直接继承 \\
    全向接收候选区域（问题二） & 只解决距离条件，\textbf{不能}承诺定向接收 \\
    无信号排除 1000 米圆盘 & \textbf{禁用}：定向源可能就在旁边但背向 \\
    全向覆盖终止证书 & 更换为对任意朝向成立的三角网证书 \\
    联合路线、每次补测计费、完整尾部统计 & 继承其原则 \\
    \bottomrule[1.2pt]
  \end{tabularx}
  \endgroup
\end{table}

一个说明"外侧扫描为何必要"的极端例子是：源位于 $g=(1800,0)$、朝向 $n=(1,0)$。此时严格位于场地内部的点全部处于其发射半平面的后方，仅在场内扫描没有任何发现保证。因此允许机器狗走出半径 1800 米的源分布区域是覆盖逻辑的内在需要，而不是数值越界；题目附件亦明确允许场外位置。

\subsection{状态表示与更新}

机器狗的状态为当前位置 $p$、当前检测频道 $c$、已发现频道集合 $D$ 与已清除频道集合 $C$。未知频道指尚未被发现、也尚未被证明不存在的频道。历史动作按频道保存其位置与真实响应；仅经过某点不计作检测。

评价仍以虚拟总时间为目标，形式与式 \eqref{eq:time-ledger} 相同：
\begin{equation}
  T=\frac15\sum_k\lVert p_k-p_{k-1}\rVert+5M+S+5C_s+3C_f ,
  \label{eq:q4-time}
\end{equation}
其中 $M$ 为检测次数，$S$ 为检测频道切换次数，$C_s$、$C_f$ 分别为成功与失败的光学操作次数。清除操作不改变检测频道，不要求返回原点。评价同时报告清除完整率与 $T/N_{\text{clear}}$；本地全清场景以真实 $N$ 计算 $T/N$ 并对场景等权平均，未清完场景的低时间不得当作有效速度。

完整的物理可行状态可写作 $(g,R,n,\tau)$ 的集合 $\Omega$。正向观测交入接收条件与测角约束；负向观测交入
\begin{equation}
  \lVert s-g\rVert>R
  \quad\text{或}\quad
  \bigl[\tau=\text{定向},\ n^{T}(s-g)<0\bigr].
  \label{eq:neg-obs}
\end{equation}
该析取导致联合状态非凸。本文不声称精确求解 $\Omega$，而是分开保存以下两类表示。

\textbf{可靠位置外包 $P$。} 初始为 1800 米圆域的 180 边外接多边形；正向观测交入半宽 $\varepsilon=1.0051^\circ$ 的前向角域与半径 1500.000001 米圆盘的外接多边形；近距记录交入半径 5.000001 米圆盘的外包。保留 0.005$^\circ$ 舍入余量与微小数值余量。\textbf{无信号记录只记入历史，不对 $P$ 做任何圆盘删减。} 真实位置应始终包含在 $P$ 中。

\textbf{规划用有限假设 $H$。} 取 $P$ 内的顶点与内部插值位置，对每个位置采样至多 3 个接收半径与 72 个朝向，并保留全向假设。历史负观测按"超出距离或处于背面"进行筛选。相容样本中"可接收"的比例只用于动作评分，它不是经过校准的真实概率，也不穷尽连续可行状态。样本为空只表示当前离散表示不足，\textbf{不代表源不存在}；此时评分回退为 0.5 并继续保留 $P$。

近距离观测可直接检查安全清除；接收半径未知与类型未知都不影响光学操作。

\subsection{对任意朝向成立的发现证书}
\label{sec:q4-mesh}

用边长 $L=990$ 米的等边三角网覆盖场地，基向量为 $(990,0)$ 与 $(495,495\sqrt3)$。保留与半径 1800.01 米圆域相交的全部三角形，并保留其全部顶点。实现共生成 31 个点、42 个三角形，其中 18 个点位于圆域之外；0.01 米为边界数值余量，不扩大源先验。

\textbf{定理。} 任取三角形顶点 $s_1,s_2,s_3$ 及三角形内的源位置 $g$，存在非负系数 $\lambda_i$ 且 $\sum_i\lambda_i=1$ 使 $g=\sum_i\lambda_is_i$。对任意朝向 $n$，有
\begin{equation}
  \sum_{i=1}^{3}\lambda_i\,n^{T}(s_i-g)
  =n^{T}\Bigl(\sum_i\lambda_is_i-g\Bigr)=0 .
  \label{eq:convex-comb}
\end{equation}
因此三个内积不可能都严格小于 0，至少有一个顶点位于源的闭发射半平面内。又因三角形中任意两点距离不超过 $L=990$ 米，小于接收半径下界 1000 米，所以该顶点能够接收到该源的信号；全向源自然也满足。$\square$

\textbf{推论。} 若某频道在某个三角形的三个顶点都实际测得无信号，则该三角形内不存在此频道的源。当全部三角形都满足该条件时，整个场地被排除。

这里必须登记完全相同的网格坐标，不能用邻近但不同的检测点替代顶点，否则会破坏角度边界处的凸组合依据。

作为一个可人工核验的例子：取 $s_1=(0,0)$、$s_2=(990,0)$、$s_3=(495,\,857.3651497)$，源位于重心 $g=(495,\,285.7883832)$，三个顶点到 $g$ 的距离均约为 571.5768 米。当 $n=(0,-1)$ 时，前两个顶点的内积为 $285.7884>0$，第三个为 $-571.5768$，即至少有一个顶点可接收，且都在 1000 米以内。形式证明对任意位置与任意朝向成立，不依赖这个例子，也不依赖朝向的离散采样。

该网格给出的是保守充分条件，不是必要条件，也不是最优站点布局。

\subsection{凸包排除定理与 22 点布局}

三角网证书要求三个顶点都实际无信号。实践中可以放宽为更一般的条件。

\textbf{定理（凸包排除）。} 设 $Q$ 为一个凸子区域，$S_Q$ 为同一频道的实际无信号点集合，且 $S_Q$ 中每个点 $s$ 到 $Q$ 的每个顶点的距离都不超过 999.99 米。由欧氏距离的凸性，这保证每个 $s$ 到 $Q$ 内所有点的距离都不超过该范围。若 $Q\subseteq\operatorname{conv}(S_Q)$，则 $Q$ 中不存在任何朝向的该频道源。

\textbf{证明。} 反设 $Q$ 中存在一个定向源位于 $g$，则其全部无信号点必须满足 $n^{T}(s-g)<0$。但 $g\in Q\subseteq\operatorname{conv}(S_Q)$，故 $g$ 是这些 $s$ 的凸组合；取相同权重求和得 $n^{T}(g-g)=\sum_i\lambda_in^{T}(s_i-g)<0$，即 $0<0$，矛盾。全向源更直接由距离条件排除。半平面边界包含在可接收区域内，是推导的必要条件。$\square$

原先的 31 点三角网只是该定理在固定三角网下的一种特例。新算法允许子区域 $Q$ 的不同顶点由不同的三点组合证明，只要各组合的站点都满足"到整个 $Q$ 的距离条件"即可得到 $Q$ 属于所有这些站点凸包的结论。需注意：不能仅分别证明每个顶点可被排除，而忽略整个 $Q$ 的距离条件。

\textbf{布局构造。} 据此，本文选取：中心 1 点 $(0,0)$，半径 980 米的等角内圈 7 点，半径 1850 米的等角外圈 14 点，均从正 $x$ 方向起算，共 22 个站点，替代原先的 31 点网格。以原三角网覆盖作为分区起点，每个三角形与半径 1800 米圆的外切 360 边形相交（不能用内接多边形替代场地）；对无法直接获得凸包证据的三角形递归四分。22 点构造共得到 1334 个通过条件的凸子区域并覆盖整个连续场地；25 点备选方案得到 722 个子区域但选择集效率低于 22 点，保留为备选；19 点与 21 点的更少点候选未通过当前几何条件。

执行器预先计算三点组合对每个子区域顶点的证据，并用位集加速实际负观测组合的查询。只有\textbf{真正执行过}且返回无信号的站点才能用于结束证明；未来站点在路径规划时可暂按无信号计算预期覆盖，但不进入完成证书。发现 16 个频道后不再搜索未知频道，但仍必须清除全部已发现源。

本轮的主要收益假设是：减少站点数量，并把原先最外层 2619 米的站点移至 1850 米圈，从而同时减少移动距离与搜索检测次数。这一假设必须由同场路径与检测账目支持，不能仅凭站点数量的下降认定时间改善。

\subsection{发现—定位—清除闭环}

\begin{enumerate}[label=(\arabic*),leftmargin=2.2em]
  \item 在原点逐频道扫描；同一位置的 20 个频道仍分别计时。近距记录立即清除。
  \item 将仍需要扫描的网格站点与已发现未清除源的区域中心作为任务点，构造固定起点、自由终点的路线；最近邻初始化后最多做 30 轮 2-opt，只执行当前路线的第一项，观测后重新计算。
  \item 若任务为扫描，则检测尚未被排除、且从未在该顶点返回负观测的未知频道；默认不附加已知频道的测量。
  \item 若任务为定位，先求保证清除点，若存在则执行。否则在区域中心、长轴两侧、末次正向观测点向区域中心前进 45\% 与 75\% 的位置以及侧移点中选下一测点，排除与历史测点距离不足 8 米的位置。
  \item 连续 3 次无信号或 10 次局部观测仍未达到清除精度时，转入光学区域后备，不靠原地重测或强行删除区域来结束。
  \item 已发现 16 个不同频道后立即停止未知频道搜索，只完成已知源；清除 16 个后结束。若不足 16 个，则必须所有已知源均已清除，且每个其余频道都有全场排除证书，才能结束。
\end{enumerate}

路线中的区域中心是服务位置的近似，并未完整预测定位进入点与清除退出点，因此它不是全局在线最优路线。

\subsubsection{下一测点的时间评分}

对候选点 $q$，记 $H$ 中可接收比例估计为 $\hat p$，围绕预测中心方位的 $-1^\circ$、$0^\circ$、$+1^\circ$ 三个假想读数给出平均区域半径 $\hat r'$。它们只估计测量收益，\textbf{不是}所有可能读数下的严格最坏情况。再记位置中心为 $c$、末次正向观测位置为 $s^+$、当前包围半径为 $\varrho$，评分取
\begin{equation}
  J(q)=\frac{\lVert p-q\rVert}{5}+6
  +\hat p\Bigl(\frac{\lVert q-c\rVert}{5}+\frac{2\hat r'}{5}\Bigr)
  +(1-\hat p)\Bigl(\frac{\lVert q-s^+\rVert}{5}+\frac{2\varrho}{5}+20\Bigr).
  \label{eq:q4-score}
\end{equation}
选择分数最小的点。式中距离系数 $1/5$ 来自移动速度；固定项 6 秒近似一次检测与换频（同频道实际仅 5 秒，但该公共项不改变同频道候选之间的排序）；$2\varrho$ 用于近似剩余定位行程；20 秒为失联处理的工程惩罚，不是题目给定参数，也不是对真值的估计。侧移距离 $\max(30,\min(180,0.3\varrho))$ 同样是工程参数。所有真正执行的动作费用按接口重新记账。该评分把第三问"少走路却多测、抵消收益"的经验用于决策，但当前只实现了局部评分，其权重仍需通过后续消融与官方演练检验。

\subsection{清除与光学后备}

常规清除采用 $\mathcal{K}=\bigcap_{v\in\text{vertices}(P)}\overline{B}(v,19.5)$。若 $\mathcal{K}$ 非空，选择距机器狗最近的点 $q\in\mathcal{K}$；对凸多边形，距 $q$ 的最大值在顶点取得，故真实源到 $q$ 的距离不超过 $19.5<20$ 米。已认证清除若执行失败，则记录冲突并停止，不能继续宣称保证。

光学后备搜索把 $P$ 投影到其直径方向的坐标系，以边长 26 米的正方格覆盖，保留所有与 $P$ 相交的方格，在每个方格中心尝试一次清除。每格内任一点到格心的距离不超过 $13\sqrt2\approx18.3848$ 米，故当全部方格被覆盖时必有一次成功，且该结论与源的发射朝向无关。实际采用最近邻次序，成功即停；失败的尝试各支付 3 秒及相应移动费用。它是有限覆盖搜索，不是每一次尝试都保证成功，也不被伪装成零失败清除。

上述完整性逻辑在题面静态物理条件成立、观测误差界成立、动作能正常执行且未被现实或虚拟预算截断时有效。数值几何采用浮点与外包余量，并非形式化区间算术证明；预算中止或异常必须作为未完成状态保存。

\subsection{求解结果}

在 40 场冻结留出配对中（种子 26000--26039），原 31 点网格策略与 22 点布局策略使用完全相同的场景。

\begin{table}[H]
  \centering
  \caption{问题四 40 场冻结配对结果（场景等权 $T/N$，单位：秒/点）}
  \label{tab:q4-main}
  \begingroup
  \renewcommand{\arraystretch}{1.4}
  \begin{tabularx}{\textwidth}{L{0.30\textwidth}C{0.14\textwidth}C{0.14\textwidth}C{0.14\textwidth}Y}
    \toprule[1.2pt]
    {\heifont\bfseries 策略} & {\heifont\bfseries 时间/秒} & {\heifont\bfseries 平均路程/米} & {\heifont\bfseries 平均检测次数} & {\heifont\bfseries 说明} \\
    \midrule[0.5pt]
    31 点网格（原版） & 696.78 & 36537.13 & 312.43 & 全部完整清除 \\
    22 点布局（本模型） & 515.63 & 26611.70 & 249.95 & 全部完整清除 \\
    \midrule[0.5pt]
    平均节省 & 181.15 & 9925.43 & 62.48 & 降低 26.00\% \\
    \bottomrule[1.2pt]
  \end{tabularx}
  \endgroup
\end{table}

40 场平均由 696.78 秒/点降至 515.63 秒/点，节省 181.15 秒/点，降低 26.00\%；配对节省的近似 95\% 区间为 $[151.78,210.53]$ 秒/点。40 场中 39 场改善、1 场退化。平均路程由 36537.13 米降至 26611.70 米，平均检测次数由 312.43 次降至 249.95 次。

\textbf{必须报告退化案例。} 最差场景（种子 26011）反而增加 145.79 秒/点；该场的源数 $N=16$，原策略较早发现了全部频道而提前停止搜索，而 22 点布局的保底覆盖站点较少，在该场次中未能同样早地触发停止。这说明\textbf{更少的保底覆盖站点并不保证每个场景都更快}，本节的开头结论只应表述为"平均改善"，而不是"任意场景均更快"。

此外还进行了 48 场压力与类型比例测试，结果见表 \ref{tab:q4-stress}。这些测试与主 40 场分开统计，压力案例的均值不混入主均值。

\begin{table}[H]
  \centering
  \caption{问题四压力与类型比例测试（各组均完整清除，单位：秒/点）}
  \label{tab:q4-stress}
  \begingroup
  \renewcommand{\arraystretch}{1.4}
  \begin{tabularx}{\textwidth}{L{0.30\textwidth}C{0.10\textwidth}C{0.15\textwidth}C{0.15\textwidth}C{0.15\textwidth}}
    \toprule[1.2pt]
    {\heifont\bfseries 测试组} & {\heifont\bfseries 场数} & {\heifont\bfseries 原版} & {\heifont\bfseries 22 点版} & {\heifont\bfseries 平均节省} \\
    \midrule[0.5pt]
    边界外向 / 恒定 $+1^\circ$   & 8  & 950.49  & 743.54 & 206.95 \\
    边界切向 / 恒定 $-1^\circ$   & 8  & 745.68  & 522.47 & 223.22 \\
    近源聚集 / 棋盘式误差场      & 8  & 1041.71 & 800.07 & 241.65 \\
    $N-1$ 个定向源               & 12 & 817.19  & 603.25 & 213.94 \\
    仅 1 个定向源                & 12 & 638.18  & 432.14 & 206.04 \\
    \bottomrule[1.2pt]
  \end{tabularx}
  \endgroup
\end{table}

48 场压力与类型比例测试全部完整清除，且各组均值均有改善。需要说明：源生成分布仍是本地设定，压力组的均值不混入主均值；覆盖构造的数值核验不等于策略全局最优，也不构成对任意布局的效率保证。

\subsection{有效性检验与适用边界}

\textbf{独立覆盖核验。} 规划器内部使用三角形叉积与位集判定；独立审核则从日志中提取真实无信号点，用凸包半空间表示检查每一个子区域，另外进行分区并集与场地差集检查，并重算每个动作的时间。主审计与压力组共 88 场、28497 个动作、701684 个子区域检查通过独立审核。几何实现为浮点数加明确容差，不是区间算术证明。

\textbf{方案取舍记录。} 仅在 31 点上改进证明、增加额外清除点全频道扫描、以目标位置报价等尝试均未获得收益，予以保留但不采用；19 点与 21 点的更少点候选未通过当前几何条件，\textbf{不能}据此证明所有更少点布局都不可能；25 点通过几何条件但选择集效率低于 22 点，保留为备选。本轮采用的 22 点布局运行保留配置、场景、源码快照与散列以及逐动作记录，新版单场最大计算时间 0.93 秒；这些记录不以失败光学尝试为理由删场。

\textbf{适用边界。} 覆盖构造的数值核验不等于策略的全局最优，也不构成对任意布局的效率保证。源位置、类型与朝向的生成分布未知，题目未给定定向源的比例先验，本文只报告本地构造场景下的结果。此外，策略的有效性依赖于题面静态物理条件与观测误差界成立；若接口的实际行为与题述不符，完整性论证不再适用。

\subsection{正式测试结果}

与其他要求相同，第四问需在充分演练后进行三次正式测试并按表 1 格式给出结果。

\begin{table}[H]
  \centering
  \caption{问题四正式测试结果}
  \label{tab:q4-formal}
  \begingroup
  \renewcommand{\arraystretch}{1.45}
  \begin{tabularx}{\textwidth}{L{0.22\textwidth}C{0.18\textwidth}C{0.20\textwidth}C{0.20\textwidth}}
    \toprule[1.2pt]
    {\heifont\bfseries 测试案例编码} & {\heifont\bfseries 清除干扰源个数} & {\heifont\bfseries 平均定位清除时间} & {\heifont\bfseries 程序运行时间} \\
    \midrule[0.5pt]
    测试1 \hspace{2em} & \hspace{2em} & \hspace{2em} & \hspace{2em} \\
    测试2 \hspace{2em} & \hspace{2em} & \hspace{2em} & \hspace{2em} \\
    测试3 \hspace{2em} & \hspace{2em} & \hspace{2em} & \hspace{2em} \\
    \bottomrule[1.2pt]
  \end{tabularx}
  \endgroup
\end{table}

\textbf{截至本文写作时，第四问的三次正式测试尚未进行，表 \ref{tab:q4-formal} 的内容留空。} 第四问的全部数值均来自本地合成场景，尚无官方演练或正式测试记录，相应的三份模拟器日志文件将在测试完成后导出并放入支撑材料。
